% GENERATED FILE. Edit canonical lessons, then run npm run generate:books.
% canonical-source-hash: fnv1a64:18bd8011baf3ee44
% canonical-lessons: HI-C356-barabar, HI-C356-bhedbhav, HI-C356-paksapat, HI-C356-imsaf, HI-C356-uksana

\chapter{Fair and Unfair}
\label{ch:hi-fair-and-unfair}
\hlchaptermodality{356}

\begin{chapteropening}
\textbf{By the end of this chapter:} \emph{I can say whether people are treated fairly.}

Complete the last lesson of chapter 356: I can say whether people are treated fairly.
\end{chapteropening}

\section[\texorpdfstring{\hi{बराबर} (barābar)}{बराबर (barābar)}]{\hi{बराबर} (barābar) — equal}

\label{lesson:HI-C356-barabar}



\begin{quote}
Before the new one: say the Hindi for a command, then the Hindi for to suggest.
\end{quote}

\subsection*{You'll want to know: \hi{बराबर}}
\textbf{\hi{बराबर}} — \emph{barābar} — \textquotedblleft{}equal\textquotedblright{}.

\begin{quote}
\hi{क़ानून} \hi{के} \hi{सामने} \hi{सब} \hi{लोग} \hi{बराबर} \hi{हैं।}
\end{quote}

\begin{quote}
\emph{qānūn ke sāmne sab log barābar haiṁ.}
\end{quote}

\begin{quote}
Everyone is equal before the law.
\end{quote}

\begin{cousinweb}
\textbf{\hi{बराबर}} (\emph{barābar}) is usually traced to Persian \emph{barābar}, \textquotedblleft{}equal, level\textquotedblright{}.
\end{cousinweb}

\begin{grammarlens}[title={What you've built}]
One more word for giving reasons and opinions.
\end{grammarlens}

\subsection*{Your turn}
\begin{itemize}
\raggedright
  \item \emph{Say it:} \emph{barābar}
  \item \emph{Say it:} \emph{barābar}, once more
  \item \emph{Say it:} \emph{qānūn ke sāmne sab log barābar haiṁ.} — Everyone is equal before the law.
  \item \emph{Recall:} say the Hindi for a command, then the Hindi for to suggest, then say \emph{barābar} again
\end{itemize}

\begin{tcolorbox}[breakable,colback=teal!4,colframe=teal!35!black,title={Before you move on}]
What does \hi{बराबर} mean? (\textquotedblleft{}Equal\textquotedblright{}.) Say it once more.
\end{tcolorbox}
\section[\texorpdfstring{\hi{भेदभाव} (bhedbhāv)}{भेदभाव (bhedbhāv)}]{\hi{भेदभाव} (bhedbhāv) — discrimination}

\label{lesson:HI-C356-bhedbhav}



\begin{quote}
Before the new one: say the Hindi for to suggest, then the Hindi for equal.
\end{quote}

\subsection*{You'll want to know: \hi{भेदभाव}}
\textbf{\hi{भेदभाव}} — \emph{bhedbhāv} — \textquotedblleft{}discrimination\textquotedblright{}.

\begin{quote}
\hi{स्कूल} \hi{में} \hi{लड़कियों} \hi{के} \hi{साथ} \hi{भेदभाव} \hi{नहीं} \hi{होना} \hi{चाहिए।}
\end{quote}

\begin{quote}
\emph{skūl meṁ laṛkiyoṁ ke sāth bhedbhāv nahīṃ honā cāhiye.}
\end{quote}

\begin{quote}
There should be no discrimination against girls at school.
\end{quote}

Masculine.

\begin{cousinweb}
\textbf{\hi{भेदभाव}} (\emph{bhedbhāv}) joins two Sanskrit words, \emph{bheda}, \textquotedblleft{}a split, a distinction\textquotedblright{}, and \emph{bhāva}, \textquotedblleft{}a state, an attitude\textquotedblright{} (Monier-Williams): treating people as split apart.
\end{cousinweb}

\begin{grammarlens}[title={What you've built}]
One more word for giving reasons and opinions.
\end{grammarlens}

\subsection*{Your turn}
\begin{itemize}
\raggedright
  \item \emph{Say it:} \emph{bhedbhāv}
  \item \emph{Say it:} \emph{bhedbhāv}, once more
  \item \emph{Say it:} \emph{skūl meṁ laṛkiyoṁ ke sāth bhedbhāv nahīṃ honā cāhiye.} — There should be no discrimination against girls at school.
  \item \emph{Recall:} say the Hindi for to suggest, then the Hindi for equal, then say \emph{bhedbhāv} again
\end{itemize}

\begin{tcolorbox}[breakable,colback=teal!4,colframe=teal!35!black,title={Before you move on}]
What does \hi{भेदभाव} mean? (\textquotedblleft{}Discrimination\textquotedblright{}.) Say it once more.
\end{tcolorbox}
\section[\texorpdfstring{\hi{पक्षपात} (pakṣapāt)}{पक्षपात (pakṣapāt)}]{\hi{पक्षपात} (pakṣapāt) — bias, favouritism}

\label{lesson:HI-C356-paksapat}



\begin{quote}
Before the new one: say the Hindi for equal, then the Hindi for discrimination.
\end{quote}

\subsection*{You'll want to know: \hi{पक्षपात}}
\textbf{\hi{पक्षपात}} — \emph{pakṣapāt} — \textquotedblleft{}bias, favouritism\textquotedblright{}.

\begin{quote}
\hi{जज} \hi{को} \hi{पक्षपात} \hi{नहीं} \hi{करना} \hi{चाहिए।}
\end{quote}

\begin{quote}
\emph{jaj ko pakṣapāt nahīṃ karnā cāhiye.}
\end{quote}

\begin{quote}
A judge should not show bias.
\end{quote}

Masculine.

\begin{cousinweb}
\textbf{\hi{पक्षपात}} (\emph{pakṣapāt}) is a tatsama, borrowed from Sanskrit \emph{pakṣapāta}, \textquotedblleft{}taking a side, partiality\textquotedblright{} (Monier-Williams).
\end{cousinweb}

\begin{grammarlens}[title={What you've built}]
One more word for giving reasons and opinions.
\end{grammarlens}

\subsection*{Your turn}
\begin{itemize}
\raggedright
  \item \emph{Say it:} \emph{pakṣapāt}
  \item \emph{Say it:} \emph{pakṣapāt}, once more
  \item \emph{Say it:} \emph{jaj ko pakṣapāt nahīṃ karnā cāhiye.} — A judge should not show bias.
  \item \emph{Recall:} say the Hindi for equal, then the Hindi for discrimination, then say \emph{pakṣapāt} again
\end{itemize}

\begin{tcolorbox}[breakable,colback=teal!4,colframe=teal!35!black,title={Before you move on}]
What does \hi{पक्षपात} mean? (\textquotedblleft{}Bias, favouritism\textquotedblright{}.) Say it once more.
\end{tcolorbox}
\section[\texorpdfstring{\hi{इंसाफ़} (iṁsāf)}{इंसाफ़ (iṁsāf)}]{\hi{इंसाफ़} (iṁsāf) — fairness, justice (everyday)}

\label{lesson:HI-C356-imsaf}



\begin{quote}
Before the new one: say the Hindi for discrimination, then the Hindi for bias.
\end{quote}

\subsection*{You'll want to know: \hi{इंसाफ़}}
\textbf{\hi{इंसाफ़}} — \emph{iṁsāf} — \textquotedblleft{}fairness, justice (everyday)\textquotedblright{}.

\begin{quote}
\hi{यह} \hi{इंसाफ़} \hi{नहीं} \hi{है}!
\end{quote}

\begin{quote}
\emph{yah iṁsāf nahīṃ hai!}
\end{quote}

\begin{quote}
This isn't fair!
\end{quote}

Masculine.

\begin{cousinweb}
\textbf{\hi{इंसाफ़}} (\emph{iṁsāf}) is usually traced to Arabic \emph{inṣāf}, \textquotedblleft{}equity\textquotedblright{}, which reached Hindi through Persian.
\end{cousinweb}

\begin{grammarlens}[title={What you've built}]
One more word for giving reasons and opinions.
\end{grammarlens}

\subsection*{Your turn}
\begin{itemize}
\raggedright
  \item \emph{Say it:} \emph{iṁsāf}
  \item \emph{Say it:} \emph{iṁsāf}, once more
  \item \emph{Say it:} \emph{yah iṁsāf nahīṃ hai!} — This isn't fair!
  \item \emph{Recall:} say the Hindi for discrimination, then the Hindi for bias, then say \emph{iṁsāf} again
\end{itemize}

\begin{tcolorbox}[breakable,colback=teal!4,colframe=teal!35!black,title={Before you move on}]
What does \hi{इंसाफ़} mean? (\textquotedblleft{}Fairness, justice (everyday)\textquotedblright{}.) Say it once more.
\end{tcolorbox}
\section[\texorpdfstring{\hi{उकसाना} (uksānā)}{उकसाना (uksānā)}]{\hi{उकसाना} (uksānā) — to provoke, to stir up}

\label{lesson:HI-C356-uksana}



\begin{quote}
Before the new one: say the Hindi for bias, then the Hindi for fairness.
\end{quote}

\subsection*{You'll want to know: \hi{उकसाना}}
\textbf{\hi{उकसाना}} — \emph{uksānā} — \textquotedblleft{}to provoke, to stir up\textquotedblright{}.

\begin{quote}
\hi{कुछ} \hi{लोग} \hi{भीड़} \hi{को} \hi{उकसा} \hi{रहे} \hi{थे।}
\end{quote}

\begin{quote}
\emph{kuch log bhīṛ ko uksā rahe the.}
\end{quote}

\begin{quote}
Some people were stirring up the crowd.
\end{quote}

Past, with \hi{ने} (he or she …): \hi{उसने} \hi{भीड़} \hi{को} \hi{उकसाया}, \emph{usne bhīṛ ko uksāyā}.

\begin{cousinweb}
\textbf{\hi{उकसाना}} (\emph{uksānā}) is a tadbhava, inherited from Sanskrit \emph{utkasati}, \textquotedblleft{}opens\textquotedblright{} (CDIAL 1718). CDIAL glosses the Hindi verb \textquotedblleft{}to draw out, to raise, to excite\textquotedblright{}.
\end{cousinweb}

\begin{grammarlens}[title={What you've built}]
One more word for giving reasons and opinions.
\end{grammarlens}

\subsection*{Your turn}
\begin{itemize}
\raggedright
  \item \emph{Say it:} \emph{uksānā}
  \item \emph{Say it:} \emph{uksānā}, once more
  \item \emph{Say it:} \emph{kuch log bhīṛ ko uksā rahe the.} — Some people were stirring up the crowd.
  \item \emph{Recall:} say the Hindi for bias, then the Hindi for fairness, then say \emph{uksānā} again
\end{itemize}

\begin{tcolorbox}[breakable,colback=teal!4,colframe=teal!35!black,title={Before you move on}]
What does \hi{उकसाना} mean? (\textquotedblleft{}To provoke, to stir up\textquotedblright{}.) Say it once more.
\end{tcolorbox}
